\section{Introduction}

\subsection{Context of Underwater Robotics and ROVs}
\begin{itemize}
	\item Overview of underwater robotic platforms: ROVs vs AUVs
	\item Typical ROV applications: inspection, maintenance, research, search and recovery
	\item Operational constraints: communications, autonomy limits, energy management
\end{itemize}

\subsection{Why Docking Is Important}
\begin{itemize}
	\item Enables in-water recharge and battery swaps to extend mission duration
	\item Allows data offload without surfacing, maintaining mission continuity
	\item Reduces manual recovery, vessel time and associated operational risks/costs
	\item Supports persistent monitoring and multi-mission deployments
\end{itemize}

\subsection{Limitations of Current ROV Operations}
\begin{itemize}
	\item Dependence on tethers for power and communications limits range and mobility
	\item Finite battery life restricts mission duration and payloads
	\item Manual recovery and surface operations are weather and resource dependent
	\item Existing docking research often targets AUVs or static, seabed-mounted docks
\end{itemize}

\subsection{What a Free-Floating Docking Station Is}
\begin{itemize}
	\item Definition: an untethered, buoyant station that maintains position or drifts
	\item Advantages: relocatability, compatibility with mobile operations, reduced seabed impact
	\item New challenges: dock platform motion in surge/sway/heave and rotational degrees of freedom
\end{itemize}

\subsection{Why Vision-Based Docking Is Difficult}
\begin{itemize}
	\item Turbidity: scattering and absorption reduce contrast and visibility
	\item Variable lighting: low light, shadows, sun glint and inhomogeneous illumination
	\item Relative motion from currents and dock movement complicates tracking and control
	\item Image distortions: refraction, lens distortion and underwater optical effects
	\item Onboard compute limits require real-time, efficient algorithms
\end{itemize}

\subsection{High-Level Description of the Proposed System}
\begin{itemize}
	\item Turbidity-adaptive vision pipeline for marker detection and image enhancement
	\item Pose estimation module producing relative 6-DOF dock-to-vehicle transform
	\item Motion compensation and control loop for approach and engagement
	\item System-level behaviours: search, approach, align, engage, abort/fallback
\end{itemize}

\subsection{Thesis Structure}
\begin{itemize}
	\item Summary of each chapter/section and how they connect (literature, methods, experiments)
	\item Outline of experimental validation: tank trials, metrics and evaluation criteria
	\item Brief note on expected results and where to find them in the report
\end{itemize}
